\chapter*{Discoveries at a Glance}
\addcontentsline{toc}{chapter}{Discoveries at a Glance}
\markboth{Discoveries at a Glance}{Discoveries at a Glance}

% Sources: CL-001--CL-006, CL-009, CL-011--CL-024; CONTENT-LEDGER BG-U-01--14;
% authoritative conditional multi-model manifest and current LSE run.
\begin{tcolorbox}[
  title={Discovery 1: current market reference and primary conditional signal},
  colback=sflightblue,
  colframe=sfblue,
  boxrule=0.7pt,
  arc=1mm,
  left=3mm,
  right=3mm
]
\small
The latest completed LSE close available to the study is \textbf{USD 47.27 per
\texttt{B} share on 27 August 2026}, verified in the 28 August extraction. In
the primary Full Bates--Hawkes scenario, the conditional Barrick distribution
has a \textbf{median of USD 34.94 per share}. The observed close is therefore
\textbf{USD 12.33, or 35.27\%, above the model median}; it lies at the
\textbf{62.27th percentile}, with \textbf{37.73\% of simulated values above
it}.

\medskip
\textbf{Directional reading.} If the median is used as the central benchmark,
the observed close is \textbf{above this conditional model median}. This does
not establish market mispricing:
USD 47.27 remains inside a wide scenario distribution, while the commodity
bridge, accounting reconciliation, equity bridge and diluted-share input are
not yet sufficient for a fair-value claim. The comparison is not a target
price and does not constitute investment advice.
\end{tcolorbox}

The two numbers answer different questions. USD 47.27 is the freshest verified
completed close in the versioned Barrick dataset, not a live quotation. USD 34.94
is the centre of a model-conditional distribution. The market observation lies
between P50 and P75: the central scenario is below the close, but a material
right tail remains above it. The sign of the gap answers the narrow
directional gap only relative to this specification;
the width of the distribution describes how strongly the unresolved
assumptions qualify that answer.

\noindent\textbf{Why Full Bates--Hawkes is shown here.} The model is the primary
structural scenario because it combines stochastic variance with
self-exciting jumps and delivers the lowest in-sample IV RMSE on the current
27 August GLD surface: 32.49 basis points, against 56.55 for Bates--Poisson,
86.28 for Heston and 114.04 for Black--Scholes. Its branching ratio equals the
0.02 lower bound, so the current surface weakly identifies self-excitation.
``Primary'' does not mean
predictively dominant. In Team 8's one-step-ahead study, Black--Scholes has a
slightly lower pooled OOS RMSE, and none of the four candidates beats the
previous-day implied-volatility surface. For that reason, Discoveries shows one
Full Bates--Hawkes distribution, while Chapter 12 reports all four on equal
terms.

\noindent\textbf{What the four models do---and do not do.} Black--Scholes/GBM,
Heston, Bates--Poisson and Full Bates--Hawkes govern only the \emph{gold-price
layer}. They do not generate mine production or costs. Those inputs come from
Team~4's separate operating work. Barrick actuals replace 2026Q1--Q2, while
Team~4 forecasts begin at 2026Q3; this mixed operating contract is identical
in every valuation. The unified DCF, stochastic WACC and equity bridge are
common as well. Team~4's legacy NSS, implied-volatility, GBM and illustrative
EBITDA artifacts are excluded. A separate NSS curve fitted to the current LSE
Treasury cross-section is used consistently by Team~8 and the valuation.

% Generated input from the authoritative accepted run.
\IfFileExists{chapters/generated/current-valuation-summary.tex}{%
  \begin{table}[H]
\centering
\caption{Primary structural scenario: Full Bates--Hawkes conditional Barrick value quantiles.}
\label{tab:primary-bates-hawkes-valuation-quantiles}
\small
\setlength{\tabcolsep}{6pt}
\renewcommand{\arraystretch}{1.08}
\begin{tabular}{rrrr}
\toprule
\textbf{Quantile} & \textbf{EV proxy} & \textbf{Equity proxy} & \textbf{Value/share} \\
& \textbf{USD mn} & \textbf{USD mn} & \textbf{USD} \\
\midrule
P01 & -908 & -908 & -0.54 \\
P05 & 2,504 & 2,504 & 1.49 \\
P10 & 5,540 & 5,540 & 3.31 \\
P25 & 25,392 & 25,392 & 15.15 \\
P50 & 58,550 & 58,550 & 34.94 \\
P75 & 109,095 & 109,095 & 65.11 \\
P90 & 175,614 & 175,614 & 104.81 \\
P95 & 228,710 & 228,710 & 136.50 \\
P99 & 365,438 & 365,438 & 218.11 \\
\bottomrule
\end{tabular}

\par\smallskip
\footnotesize The USD~47.27 close on 27 August 2026 is USD~12.33, or 35.27\%,
above the conditional median and lies at the 62.27th percentile; 37.73\% of
paths exceed it. EV and equity coincide only because the corporate bridge
still contains explicit zero placeholders. This is a provisional research
sensitivity, not a fair value or target price.
\end{table}

\begin{table}[H]
\centering
\caption{Inputs that determine how the primary distribution must be read.}
\label{tab:conditional-multimodel-assumptions}
\scriptsize
\renewcommand{\arraystretch}{1.10}
\begin{tabularx}{\textwidth}{@{}p{2.65cm}X X@{}}
\toprule
\textbf{Block} & \textbf{Run input} & \textbf{Interpretation limit} \\
\midrule
Gold-price layer & Team~8 Full Bates--Hawkes recalibrated to 27 August GLD
calls under $\mathbb Q$; paths share a USD~4,417/oz Barrick Q2 realised-price
anchor. & Transfers distributional shape only; the anchor is a quarterly
average, not spot or a physical forecast. \\
Production and costs & Barrick actuals in 2026Q1--Q2, followed by Team~4
forecasts from 2026Q3 through 2030. & Actual and forecast quarters are marked
separately; Team~4 price and illustrative valuation artifacts are excluded. \\
Cash flow and DCF & Unified tax, growth, ROIC, WACC and perpetual-growth
methodology. & Identical across models; component margin and FCFF remain
unreconciled proxies. \\
Equity bridge & Debt, cash, minorities, non-operating assets and other claims
set to zero; 1,675.51 million-share proxy. & Neither bridge nor share count is
filing-reconciled. \\
Simulation & 8,192 paths, 260 fine steps and 20 quarters; aligned engine seed
20260827 and one common WACC array. & Only the price engine changes across
branches. \\
Date alignment & Barrick anchor quarter ends 30 June; Treasury curve 24 July;
GLD surface and \texttt{B} close 27 August; valuation run 27 August 2026.
& Dates are serialized but not fully synchronous. \\
\bottomrule
\end{tabularx}
\end{table}
%
}{%
  \gapbox{The accepted Full Bates--Hawkes summary is unavailable. No
  replacement number is supplied editorially.}%
}

% Figure copied byte-identically from the authoritative accepted run.
\IfFileExists{img/valuation/full_bates_hawkes_value_distribution.png}{%
  \begin{figure}[H]
    \centering
    \includegraphics[width=0.88\textwidth]{img/valuation/full_bates_hawkes_value_distribution.png}
    \caption{Conditional Barrick value-per-share distribution under the
    primary Full Bates--Hawkes gold-price scenario. The vertical market
    reference is the completed \texttt{B} close of USD 47.27 on 27 August
    2026. Across 8,192 paths, P10/P50/P90 are USD 3.31/34.94/104.81 and the
    close lies at the 62.27th percentile. Operations combine Barrick actuals
    with explicitly labelled Team~4 forecasts; this is not a target-price
    distribution.}
    \label{fig:primary-bates-hawkes-value-distribution}
  \end{figure}
}{%
  \gapbox{The primary Full Bates--Hawkes distribution is unavailable until it
  is copied from the same accepted manifest as the summary table.}%
}

\noindent\textbf{How the evidence was assembled.} The front-page comparison
is the endpoint of an integrated chain rather than the output of an isolated
gold model. Team 8 supplies option-implied gold-price dynamics; Team 4 supplies
the ad hoc production and unit-cost layer; Team 5 maps the resulting component
margin through cash-flow, discounting and per-share identities; Team 3
supplies simulation diagnostics. The remaining teams provide probability,
econometric, historical-volatility and stochastic-process foundations. Their
unification required explicit ownership of each variable, common units and
dates, deterministic handoff contracts, source and claim freezes, and separate
code, factual, editorial and visual-review gates.

That chain also separates calibration from forecasting. Team 8 calibrates GLD
calls under the risk-neutral measure. The conditional run restarts every model
at Barrick's Q2 2026 average realised price of USD~4,417/oz and transfers only
the option-implied distributional shape. It performs neither a one-for-one GLD
level conversion nor a validated $\mathbb Q$-to-$\mathbb P$ transformation.
Moreover, its inputs are asynchronous: the anchor quarter ends 30 June, the
  Treasury curve is dated 24 July, and the GLD calibration surface and Barrick
  close are dated 27 August 2026. These choices make the experiment useful
for controlled model comparison, not a synchronous physical-gold forecast.

\begin{table}[H]
\centering
\caption{Evidence supporting the conditional valuation layer.}
\label{tab:discoveries-evidence}
\small
\renewcommand{\arraystretch}{1.13}
\begin{tabularx}{\textwidth}{@{}p{2.65cm}X X@{}}
\toprule
\textbf{Evidence block} & \textbf{Accepted current observation} &
\textbf{Permitted interpretation} \\
\midrule
Barrick equity, \texttt{B} & 86 valid daily unadjusted-close bars from
27 April to 27 August 2026; cumulative change 16.20\%, annualised log-return
volatility 51.23\%, maximum drawdown $-26.10\%$. & Short-window market
diagnostic only; not a long-run return or volatility estimate. \\
GLD reference series & 1,425 daily bars from 4 January 2021 to 27 August 2026;
total return 131.06\%, annualised volatility 18.21\%, maximum drawdown
$-26.23\%$. & ETF price in USD/share; neither physical gold in USD/oz nor
Barrick's realised selling price. \\
Current GLD surface & 24 admitted call nodes across three maturities and 13
strikes at the 27 August calibration; Treasury curve 24 July. & Current
option-implied calibration input; IV-mapped prices do not identify physical
drift. \\
Model engineering & Team 8 source freeze unchanged, 56/56 source tests and a
61-strike deterministic benchmark with maximum absolute error
$1.1611\times10^{-8}$. & Unit parity for the frozen implementation; not
historical market parity or predictive dominance. \\
Integrated experiment & 8,192 paths; four gold engines; identical separated
operating contract, unified DCF and WACC shocks; branch suite 54/54 PASS and
figure-provenance audit 47/47 PASS. & A controlled attribution of Barrick distributional differences
to the gold-price engine, subject to the declared bridge limitations. \\
\bottomrule
\end{tabularx}
\end{table}

\noindent\textbf{Who produced the study.} Research is the Starting Finance
Club PoliTo area in which members collaborate on quantitative, economic,
financial and methodological projects. The Barrick programme joined eight
autonomous teams with backgrounds spanning management engineering,
mathematics, data science, computation and quantitative finance. Knowledge
sharing was part of the research design: common foundations and coordination
rules allowed specialised work to remain deep while still being reviewed and
reused across teams. The unified article preserves their authorship and source
evidence, while treating integration as a distinct phase with its own code,
writing, technical-review, editorial-review and visual-QA responsibilities.

The chapters that follow reconstruct the company and data perimeter, the
mathematical and empirical model layers, the common operating and DCF
contracts, the four-model experiment and the validation boundaries that keep a
conditional signal from becoming an unsupported investment conclusion.
